Bei konstanter Fallgeschwindigkeit ist die resultierende Kraft null, die Luftwiderstandskraft ist gleich gross wie die Gewichtskraft:
\begin{align}
 \frac{1}{2}\,c_w\,\rho\,A\,v^2 &= m\,g
\end{align}
Die quer fallende Rolle hat die Stirnfläche $A = l\cdot d$. Für einen quer angeströmten Kreiszylinder mit $l/d = \ldRP{0}{2}$ ist (Tabellenwert für $l/d = 2$) $c_w = \cwRO$ (FoTaBe S.~189), die Dichte der Luft ist $\rho = \rhLO$ (FoTaBe S.~187). Aufgelöst nach $v$:
\begin{align}
 v &= \vRoF \\
   &= \sqrt{\frac{2\cdot\mRo\cdot\gE}{\cwR\cdot\rhL\cdot\lRo\cdot\dRo}} \\
   &= \vRo \approx \result{\vRoP{0}{2}}
\end{align}
Zur Kontrolle lässt man die Rolle aus einigen Metern Höhe fallen und misst die Fallzeit für die letzten Meter. Die konstante Geschwindigkeit wird nach kurzer Zeit erreicht; Abweichungen entstehen u.a., weil die Rolle sich beim Fallen dreht und an den Enden offen ist.
